% META: {"id": "TCOMPL-CORR-EX-007", "chapitre": "TCOMPL-CORRELATION-CAUSALITE", "type_objet": "exercice", "capacites": ["TCOMPL-CORR-C2"], "capacites_codes": ["C2"], "methodes": ["M2"], "parcours": 2, "competences": ["calculer"], "duree_min": 15, "mode_creation": "ex_nihilo", "sources_inspiration": [], "fichier_tex": "chapitres/TCOMPL-CORRELATION-CAUSALITE/exercices/TCOMPL-CORR-EX-007.tex", "corrige_tex": "chapitres/TCOMPL-CORRELATION-CAUSALITE/corriges/TCOMPL-CORR-CO-007.tex", "status": "generated"}
% BEGIN-VERIFY
% from sympy import Rational, sqrt, N
% xs = [1, 2, 3, 4, 5]
% ys = [4, 7, 9, 12, 13]
% n = len(xs)
% xbar, ybar = Rational(sum(xs), n), Rational(sum(ys), n)
% cov = Rational(1, n) * sum((x - xbar) * (y - ybar) for x, y in zip(xs, ys))
% vx = Rational(1, n) * sum((x - xbar) ** 2 for x in xs)
% vy = Rational(1, n) * sum((y - ybar) ** 2 for y in ys)
% assert cov == Rational(23, 5) and vx == 2 and vy == Rational(54, 5)
% a = cov / vx
% b = ybar - a * xbar
% assert a == Rational(23, 10) and b == Rational(21, 10)
% assert a * xbar + b == ybar
% r = cov / sqrt(vx * vy)
% assert 0.989 < float(N(r)) < 0.991
% END-VERIFY
\begin{exercice}{TCOMPL-CORR-EX-007}{2}{15}
On reprend le nuage précédent.
\begin{enumerate}
  \item Calculer la covariance de $(x\,;\,y)$ et les variances de $x$ et
        de $y$.
  \item En déduire l'équation de la droite de régression de $y$ en $x$.
  \item Calculer le coefficient de corrélation linéaire et l'interpréter.
\end{enumerate}
\end{exercice}
